\documentclass[11pt]{article}
\usepackage[margin=1in]{geometry}
\usepackage{amsmath,amssymb}

\title{Proof of Conjecture 00000003822\\
\large Diameter of a finite-type highest-weight crystal}
\author{}
\date{}

\begin{document}
\maketitle

Let \(B(\lambda)\) be the finite-type highest-weight crystal of dominant
integral weight \(\lambda\). Use the lowering arrows
\(b\to f_i(b)\) when \(f_i(b)\) is defined. We prove that the diameter
of the underlying undirected graph is \(\langle\lambda,2\rho^\vee\rangle\).

\section*{A height and two distinguished vertices}

Let \(2\rho^\vee\) be the sum of positive coroots and set
\[
   h(b)=\langle\operatorname{wt}(b),2\rho^\vee\rangle\in\mathbb Z .
\]
Each simple reflection \(s_i\) permutes the positive coroots other
than \(\alpha_i^\vee\) and negates that one. Thus
\(s_i(\rho^\vee)=\rho^\vee-\alpha_i^\vee\), which by the reflection
formula gives \(\langle\alpha_i,\rho^\vee\rangle=1\). A lowering arrow
subtracts a simple root from the weight and therefore decreases \(h\)
by exactly \(2\).

Write \(u\) for the highest vertex and \(D=h(u)=
\langle\lambda,2\rho^\vee\rangle\). Every vertex is obtained by
applying lowering operators to \(u\). The Sch\"utzenberger--Lusztig
involution \(\xi\) reverses lowering arrows, up to the permutation
of simple-root labels induced by the longest Weyl element \(w_0\).
It satisfies \(\operatorname{wt}(\xi(b))=w_0\operatorname{wt}(b)\)
and sends \(u\) to the lowest vertex \(v\)
\cite[\S\S 3, 6.1]{azenhas}. Since \(w_0^{-1}\) sends the set of
positive coroots to the set of negative coroots, invariance
of the weight--coweight pairing gives
\[
  h(\xi(b))=\langle\operatorname{wt}(b),w_0^{-1}(2\rho^\vee)\rangle
            =-h(b).
\]
In particular \(h(v)=-D\).

\section*{Distance bounds}

For any \(b\), choose a directed path \(u\to b\) and let \(r(b)\) be
its length. The height drop shows \(2r(b)=D-h(b)\), so the length
is independent of the choice of path. Applying \(\xi\) to a directed
path from \(u\) to \(\xi(b)\) and reversing it gives a directed path
from \(b\) to \(v\). Its length is \((h(b)+D)/2=D-r(b)\).
In particular the length of a directed path \(u\to v\) is \(D\), so
\(D\) is a nonnegative integer.

For arbitrary vertices \(a,b\), the undirected walk via \(u\) has
length \(r(a)+r(b)\), and the one via \(v\) has length
\(2D-r(a)-r(b)\). These lengths sum to \(2D\), so at least one
is at most \(D\). This bounds every graph distance above by \(D\).
On any undirected walk from \(u\) to \(v\), one edge changes \(h\)
by \(2\) in absolute value. The net drop is \(2D\), requiring at least
\(D\) edges. Hence \(d(u,v)=D\), which proves the diameter formula.

\section*{Lean verification}

In the accompanying Lean~4.19.0 project, \verb|HasDiameter| characterizes the
diameter by a walk-length bound for all pairs and a pair attaining that bound.

The theorem \verb|crystal_diameter_from_weights| checks the proof using the
lowering weight change, the pairing identities, generation from the highest
vertex, and the weight and arrow properties of the involution. The
application to \(B(\lambda)\) uses the crystal and root-system facts
identified above. Run \verb|lake build| in \verb|lean/| to verify the theorem
and print its axiom dependencies.

\begin{thebibliography}{1}
\bibitem{azenhas}
O.~Azenhas, M.~Tarighat Feller, and J.~Torres,
\emph{Symplectic cacti, virtualization and Berenstein--Kirillov groups},
J. Comb. Algebra, DOI \texttt{10.4171/JCA/116}.
\end{thebibliography}

\end{document}
